%%
%% spring-lecture07.tex
%% 
%% Made by Alex Nelson <pqnelson@gmail.com>
%% Login   <alex@lisp>
%% 
%% Started on  2026-04-16T10:59:59-0700
%% Last update 2026-04-16T10:59:59-0700
%% 

\lecture{}

\begin{proposition}\label{prop:2.26}
For semisimple $M$, the following are equivalent:
\begin{enumerate}
\item $M$ Is finitely-generated as an $A$-module;
\item $M=N_{1}\oplus\cdots\oplus N_{n}$ where each $N_{i}$ is simple;
\item $M$ is Artinian;
\item $M$ is Noetherian;
\item There exists some finite $m<\infty$ such that if
  $0=M_{0}\propersubset\cdots\propersubset M_{k}$ is a finite chain of strictly
  increasing submodules (of $M$), then $k\leq m$.
\end{enumerate}
\end{proposition}

\begin{proof}
Homework.
\end{proof}

\begin{definition}\label{defn:2.27}
Let $M$ be an $A$-module. We define the \define{Radical} of $M$
to be
\begin{equation}
\Rad(M):=\begin{cases}\bigcap\{N\in S(M)\mid M/N\mbox{is simple}\}&\\
M &\mbox{if above is empty}
\end{cases}
\end{equation}
\end{definition}

\begin{lemma}\label{lemma:2.28}
\begin{enumerate}
\item $\Rad(M)\in S(M)$ --- the radical of $M$ is a submodule of $M$
\item For $N\in S(M)$, if $\Rad(M/N)=0$, then $\Rad(M)\subset N$
\item $\Rad(M/\Rad(M))=0$
\item If $M$ is semisimple, then $\Rad(M)=0$.
\end{enumerate}
\end{lemma}

\begin{proposition}\label{prop:2.29}
$M$ is finitely-generated semisimple if and only if $M$ is Artinian
  and $\Rad(M)=0$.
\end{proposition}

\subsection{Structure of Semisimple Algebras}

\begin{definition}\label{defn:3.1}
Let $A$ be an algebra over $R$. We say $A$ is \define{Semisimple} (as
an algebra) if $A$ is a semisimple right $A$-module (usually denoted $A_{A}$).

\textsc{Note}: since $A=1_{A}A$ by Proposition~\ref{prop:2.26},
$A=N_{1}\oplus\cdots\oplus N_{m}$ for some minimal right ideals
$N_{1}$, \dots, $N_{m}$.
\end{definition}

\begin{proposition}\label{prop:3.2}
\begin{enumerate}
\item $A$ is semisimple if and only if $A$ is right Noetherian (as an
  $A$-module) and $\Rad(A_{A})=0$
\item If $A$ is semisimple, then every $A$-module $N$ is
  semisimple. Moreover, the simple $A$-modules are isomorphic to
  minimal right ideals of $A$ and all minimal right ideals of $A$ are
  simple $A$-modules.
\end{enumerate}
\end{proposition}

\begin{proof}
\begin{enumerate}
\item By Proposition~\ref{prop:2.29}.
\item The first half follows from the fact that any free $A$-module is
  isomorphic to the direct sum of copies of $A$, and every $A$-module
  is a $\Hom$ of free $A$-modules, and Corollary~\ref{cor:2.16}~(1)
  [implies every $P\in S(M)$ has $P$ be isomorphic to $M/P$].

The second part is by Proposition~\ref{prop:2.10}, $N$ is a simple
$A$-module iff there exists a maximal right ideal $M$ of $A$ such that
$N\iso A/M$. Since $A$ is semisimple, by
Proposition~\ref{prop:2.15}~(3) [its lattice of submodules is
  complemented], $M$ has its complement in $S(A)$, which is minimal
and isomorphic to $A/M\iso N$. [If $N'$ is complement of $M$ in $A$,
  then $A=M\oplus N'$ and moreover $N'\iso A/M\iso N$.]\qedhere
\end{enumerate}
\end{proof}

\begin{corollary}\label{cor:3.3}
If $A$ is semisimple, then every homomorphic image of $A$ is semisimple.
\end{corollary}

\begin{proof}
Let $\theta\colon A\to B$ be an algebra morphism. Without loss of
generality, $\theta$ is surjective (otherwise we could just look at
$\theta\colon A\to\theta(A)$ and call $B:=\theta(A)$). By
Proposition~\ref{prop:3.2}~(2), $B$ is semisimple as an $A$-module. So
as an $A$-module, $B$ decomposes as a direct sum of simple
$A$-modules,
\begin{equation}
B=\bigoplus_{i\in J}N_{i},
\end{equation}
where each $N_{i}$ is a simple $A$-module. Since 
\begin{equation}
\Annihilator(N_{i})\supset\Annihilator(B_{A})=\ker(\theta)
\end{equation}
and
\begin{equation}
\Annihilator(B_{A})=\{x\in A\mid Bx=B\theta(x)=0\},
\end{equation}
then by Proposition~\ref{prop:2.4} and Lemma~\ref{lemma:2.1}~(3),
there exists a $B$-module $M_{i}$ such that $N_{i}=(M_{i})_{A}$ by
Proposition~\ref{prop:2.4} and $S(N_{i})=S(M_{i})$ by Lemma~\ref{lemma:2.1}~(3).
NB: $N_{i}$ is simple as $B$-modules. Hence the claim(/) because $B$
is a semisimple algebra.
\end{proof}

\begin{lemma}\label{lemma:3.4}
Let 
\begin{equation}
A=A_{1}\times A_{2}=\{(a_{1},a_{2})\mid a_{1}\in A_{1},a_{2}\in A_{2}\}
\end{equation}
be the direct product of algebras. Let $M$ be a right ideal of $A$.
\begin{enumerate}
\item $M$ is a right ideal of $A$
\item $\End_{A_{1}}(M)=\End_{A}(M)$
\item $\hom_{A_{1}}(M,A_{1})=\hom_{A}(M,A_{1})=\hom_{A}(M,A)$
\item $S(M_{A_{1}})=S(M_{A})$
\item If $M$ is a minimal right ideal of $A_{1}$, then $M$ is a
  minimal right ideal of $A$
\item If $N$ is a minimal right ideal of $A$, then $N$ is a minimal
  right ideal of either $A_{1}$ or $A_{2}$.
\end{enumerate}
\end{lemma}

\begin{proof}
(1)--(5) omitted since $MA_{2}\subset A_{1}A_{2}=\{(a_{1},0)(0,a_{2})\}=0$.

(6) For a minimal ideal $N$ of $A$, since
\begin{equation}
NA_{i}\subset N\cap A_{i}\subset A_{i},
\end{equation}
we have either
\begin{equation}
N=NA_{i}\quad\mbox{or}\quad NA_{i}=N\cap A_{i}=0
\end{equation}
(by the minimality of $N$). Since
\begin{equation}
0\neq N=NA_{1}+NA_{2},
\end{equation}
one of the summands must vanish $NA_{i}=0$ and the other
\begin{equation}
0\neq N=NA_{i'}=N\cap A_{i'}.
\end{equation}
Then $N$ must be a minimal ideal of $A_{i'}$.
\end{proof}

\begin{corollary}\label{cor:3.5}
Let $A_{1}$, $A_{2}$ be semisimple algebras. Then their direct product $A_{1}\times A_{2}$
is a semisimple algebra.
\end{corollary}

\begin{proof}
For $A_{i}=\bigoplus N_{i,j}$ where $N_{i,j}$ are minimal right ideals
of $A_{i}$ (and these are minimal right ideals of $A$ by
Lemma~\ref{lemma:3.4}~(5)). Then
\begin{equation}
A_{1}\times A_{2}=\left(\bigoplus N_{ij}\times
0\right)\oplus\left(\bigoplus0\times N_{j,k}\right)
\end{equation}
Hence the claim.
\end{proof}

\begin{lemma}\label{lemma:3.6}
Let $N$ be a minimal right ideal of $A$, let $x\in A$. Then either
$xN=0$ or $xN\iso N$.
\end{lemma}

\begin{proof}
Apply Schur's Lemma~\ref{lemma:2.11} to $N\to xN$, $y\mapsto xy$.
\end{proof}

\begin{lemma}\label{lemma:3.7}
Let $N$ be a right ideal of $A$ such that $N^{k}=0$, let $P$ be a
simple $A$-module. Then $PN=0$ and $N\subset\Rad(A)$.
\end{lemma}

\begin{proof}
Since $P$ is simple, $PN\in S(P)$ and thus $PN=0$ or $PN=P$. But since
$N^{k}=0$, $PN=P$ is impossible since
\begin{subequations}
  \begin{align}
0 &= PN^{k}\\
&= (PN)N^{k-1}\\
&= PN^{k-1}\\
&= (PN)N^{k-2}\\
&=\cdots=P.
  \end{align}
\end{subequations}
In particular, for $P=A/M$ where $M$ is a maximal right ideal, then
\begin{equation}
(A/M)N=0.
\end{equation}
That is to say,
\begin{equation}
N\subset\bigcap\{\mbox{maximal right ideals of }A\}=\Rad(A).
\end{equation}
Hence the result.
\end{proof}

\begin{proposition}\label{prop:3.8}
For minimal right ideals $N_{1}$, $N_{2}$, the following are equivalent:
\begin{enumerate}
\item $N_{1}\iso N_{2}$ as $A$-modules
\item $N_{1}N_{2}\neq0$
\item $\exists x\in A\ldotp N_{1}=xN_{2}$
\end{enumerate}
\end{proposition}

\begin{proof}
$(1)\implies(2)$ For $\phi\colon N_{1}\to N_{2}$ being the stipulated isomorphism,
\begin{subequations}
  \begin{align}
N_{1}N_{2} &\iso\phi(N_{1}N_{2})\\
&\iso\phi(N_{1})N_{2}\\
&\iso N_{1}N_{2}
&\neq0\mbox{ by previous lemma.}
  \end{align}
\end{subequations}

$(2)\implies(3)$ If $N_{1}N_{2}\neq0$, then there exists some $x\in N_{1}$
such that $xN_{2}\neq0$. Since $N_{1}$ is simple and $0\neq xN_{2}\subset N_{1}$,
then we must have $xN_{2}=N_{1}$.

$(3)\implies(1)$ Follows from Lemma~\ref{lemma:3.6}.
\end{proof}

\begin{lemma}\label{lemma:3.9}
Let $A$ be a semisimple algebra. Suppose
\begin{equation}
A=N_{1}\oplus\cdots\oplus N_{m}
\end{equation}
where each $N_{i}$ is a minimal right ideal (of $A$). If $N$ is a
minimal right ideal of $A$, then there exists an $i$ such that $N\iso N_{i}$.
\end{lemma}

\begin{proof}
By Proposition~\ref{prop:2.15}, $N$ has a complement $M$ in $A$, so:
\begin{equation}
A=N\oplus M.
\end{equation}
By Corollary~\ref{cor:2.16}, $M$ is also semisimple as an $A$-module.
By Proposition~\ref{prop:2.19} or Theorem~\ref{thm:2.24}, there exists
an $N_{i}$ such that $N\iso N_{i}$.
\end{proof}

\begin{corollary}\label{cor:3.10}
If $A$ is a semisimple algebra, then the set of isomorphism classes of
simple $A$-modules must be finite.
\end{corollary}

\begin{proof}
Write $A=N_{1}\oplus\cdots\oplus N_{m}$ and use Lemma~\ref{lemma:3.9}.
\end{proof}

\begin{definition}
We say an algebra $A$ is \define{Simple} if $A\neq0$ and $I(A)=\{0,A\}$
is the set of \emph{two-sided} ideals of $A$.
\end{definition}